%% AUTO-GENERATED by paper/make_tables.py: DO NOT EDIT BY HAND.
%% Regenerate after every benchmark run:  python paper/make_tables.py

\begin{table}[htbp]
\caption{\label{tab:twin}The decisive control. Each quantum kernel is compared
against an SPD-manifold kernel used in the same support vector machine, on the
same covariances, with the same tuning budget, in the same reference frame:
only the metric differs. In every setting the twin is the affine-invariant
Riemannian kernel, the one classical kernel here that shares the
reference-frame quantum kernels' invariance under congruence. The
log-Euclidean kernel lacks that invariance and is shown for comparison only.
Columns give the best quantum kernel, the twin, the log-Euclidean kernel, the
range of the five quantum-minus-twin differences, and the smallest $p$
obtained by \emph{any} of the five quantum kernels against the twin.
Table~\ref{tab:seeds} repeats the PhysioNet three-qubit row under two further
partitions.}
\begin{tabular}{@{}lcccccc@{}}
\hline
Setting & $n$ & Quantum & Twin & Log-Eucl. & $\Delta$ (range over 5) & $\min p$ \\
\hline
PhysioNet, 3\,q & 104 & 0.641 & 0.630 & 0.635 & $[-0.0012, +0.0113]$ & 0.003 \\
PhysioNet, 4\,q & 104 & 0.632 & 0.634 & 0.639 & $[-0.0088, -0.0017]$ & 0.048 \\
PhysioNet, 5\,q, filter bank & 104 & 0.640 & 0.639 & 0.637 & $[-0.0051, +0.0014]$ & 0.362 \\
PhysioNet, transfer (LOSO) & 104 & 0.674 & 0.669 & 0.669 & $[+0.0028, +0.0053]$ & 0.095 \\
BCI IV-2a, 3\,q & 9 & 0.767 & 0.766 & 0.763 & $[-0.0033, +0.0016]$ & 0.426 \\
BCI IV-2a, cross-session & 9 & 0.760 & 0.755 & 0.755 & $[-0.0085, +0.0054]$ & 0.500 \\
Cho2017, 3\,q & 52 & 0.648 & 0.638 & 0.636 & $[-0.0008, +0.0100]$ & 0.001 \\
\hline
\end{tabular}
\end{table}
